\documentclass[12pt]{article}
\usepackage{amsmath}
\usepackage{amsfonts}
\usepackage{float}
\usepackage{fancyhdr}
\usepackage{graphicx}
\usepackage{setspace}
\usepackage[colorlinks=true,linkcolor=blue, citecolor=red]{hyperref}
\usepackage{url}
\usepackage[top=.75in, left=.75in, right=.75in, bottom=1in]{geometry}

% For algorithm
\usepackage{algorithm}
\usepackage{algpseudocode}

% ============ CODE ============
\usepackage{listings}
\usepackage{xcolor}
\definecolor{codegreen}{rgb}{0,0.6,0}
\definecolor{codegray}{rgb}{0.5,0.5,0.5}
\definecolor{codepurple}{rgb}{0.58,0,0.82}
\definecolor{backcolour}{rgb}{0.95,0.95,0.92}

% Styling for the code.
\lstdefinestyle{mystyle}{
    backgroundcolor=\color{backcolour},   
    commentstyle=\color{codegreen},
    keywordstyle=\color{magenta},
    numberstyle=\tiny\color{codegray},
    stringstyle=\color{codepurple},
    basicstyle=\ttfamily\footnotesize,
    breakatwhitespace=false,         
    breaklines=true,                 
    captionpos=b,                    
    keepspaces=true,                 
    numbers=left,                    
    numbersep=5pt,                  
    showspaces=false,                
    showstringspaces=false,
    showtabs=false,                  
    tabsize=2
}
\lstset{style=mystyle}

% Table packages and Unicode / Font support
\usepackage{array}
\usepackage{longtable}
\usepackage{enumitem}
\usepackage{booktabs}
\usepackage[utf8]{inputenc}
\DeclareUnicodeCharacter{2264}{\ensuremath{\le}}
\DeclareUnicodeCharacter{2194}{\ensuremath{\leftrightarrow}}
\usepackage[T1]{fontenc}
\usepackage{lmodern}
\graphicspath{{../../../../templates/}{../../prototype/}{figures/}{./}}
\setlength{\emergencystretch}{3em}
\setlist[enumerate]{label=\arabic*.,leftmargin=*}
\setlist[itemize]{leftmargin=*}

% ============ HEADER AND FOOTER ============
\setlength{\headheight}{29.43912pt}

\pagestyle{fancy}
\fancyfoot{}
\fancyfoot[R]{Page \thepage}

\setlength{\parindent}{0pt}
\renewcommand{\baselinestretch}{1.5}

\newcommand{\coursename}{Introduction of Software Engineering}
\newcommand{\reportname}{Studitopia}
\newcommand{\reporttitle}{\href{https://github.com/RoundFloor/Studitopia.git}{Studitopia}}

\newcommand{\studentname}{%
Nguyen Thanh An (24125024)\\
Nguyen Chi Bao (24125025)\\
Tran Nguyen Khanh Binh (24125026)\\
Tran Canh Anh Tuan (24125048)\\
Nguyen Le Quoc Huy (24125100)
}
\newcommand{\teachername}{Ho Tuan Thanh, Nguyen Thi Minh Tuyen}

% Header
\lhead{\reporttitle}
\rhead{
University of Science - VNUHCM\\
\coursename
}

% Footer
\newcommand{\leftfooter}{}
\lfoot{\leftfooter}

% ============ DOCUMENT ============
\begin{document}

\pagenumbering{roman}
\begin{titlepage}
\newcommand{\HRule}{\rule{\linewidth}{0.5mm}}
\centering

\textsc{\LARGE VNU-HCM}\\[0.5cm]
\textsc{\Large University of Science}\\[0.5cm]
\textsc{\large Faculty of Information Technology}\\[0.5cm]
\textsc{Computer Science}\\[0.5cm]

\HRule \\[0.4cm]
{ 
\huge{\bfseries{\reporttitle}}\\[0.5cm]
\large{\bfseries{UC-34 --- Report Inappropriate Behavior}}
}\\[0.4cm]
\HRule \\[0.5cm]

\textbf{\large Subject: \coursename}\\[0.5cm]

\begin{minipage}[t]{0.55\textwidth}
\begin{flushleft} \large
\emph{Students:}\\
\studentname
\end{flushleft}
\end{minipage}
~
\begin{minipage}[t]{0.4\textwidth}
\begin{flushright} \large
\emph{Supervising Teachers:} \\
\teachername
\end{flushright}
\end{minipage}\\[1cm]

{\large \today}\\[1cm]

\includegraphics[scale=.20]{img/hcmus-logo.png}\\[1cm] 

\vfill
\end{titlepage}	

\tableofcontents
\pagebreak

\pagenumbering{arabic}
\setcounter{page}{1}

\hypertarget{sec:uc-34-report-inappropriate-behavior}{}
\addcontentsline{toc}{section}{UC-34 --- Report Inappropriate Behavior}
\section*{UC-34 --- Report Inappropriate Behavior}
\par\medskip

\textbf{Version:} \texttt{1.0}
\par\medskip

\textbf{Performed by:} [Tran Canh Anh Tuan] | \textbf{Edited by:} [Tran Canh Anh Tuan]
\par\medskip

\begingroup
\small
\renewcommand{\baselinestretch}{1.0}\selectfont
\renewcommand{\arraystretch}{1.15}
\setlength{\tabcolsep}{6pt}
\begin{longtable}{|>{\raggedright\arraybackslash}p{0.09\textwidth}|>{\raggedright\arraybackslash}p{0.84\textwidth}|}
\hline
\textbf{Field} & \textbf{Value} \\ \hline
\endhead
Use Case ID & UC-34 \\ \hline
Name & Report Inappropriate Behavior \\ \hline
Primary Actor & User \\ \hline
Supporting Actors & Moderator \\ \hline
Feature Group & Notifications and Safety \\ \hline
Priority & Not explicitly stated \\ \hline
Related User Stories & US-40 \\ \hline
Related Requirements & AC-40.1, AC-40.2, AC-40.3, AC-40.4; FR-40.1, FR-40.2, FR-40.3, FR-40.4, FR-40.5; NFR-SEC-1, NFR-SEC-3, NFR-SEC-6, NFR-SEC-9, NFR-SEC-10, NFR-USE-6, NFR-PRIV-7 \\ \hline
Source Documents & \texttt{proposals/userstory.md} (sections 11.2-11.3); \texttt{proposals/functional\_requirement.md} (US-40); \texttt{proposals/settings.md}; \texttt{proposals/friend.md}; \texttt{proposals/friend\_profile.md}; \texttt{proposals/group\_detail.md}; \texttt{proposals/moderation\_queue.md}; \texttt{proposals/non\_functional\_requirement.md}; \texttt{proposals/API.md}; \texttt{proposals/response\_code.md}; \texttt{proposals/feature\_flow.md}; \texttt{proposals/screen\_flow.md} \\ \hline
\end{longtable}
\endgroup
\par\medskip

\hypertarget{sec:part-a-use-case-model}{}
\addcontentsline{toc}{subsection}{Part A --- Use Case Model}
\subsection*{Part A --- Use Case Model}
\par\medskip

\hypertarget{sec:a1-focused-use-case-diagram}{}
\addcontentsline{toc}{subsubsection}{A1. Focused Use Case Diagram}
\subsubsection*{A1. Focused Use Case Diagram}
\par\medskip

\begin{center}
\includegraphics[width=0.85\textwidth,height=0.75\textheight,keepaspectratio]{"figures/UC-34-report-inappropriate-behavior_1.png"}
\end{center}
\par\medskip

\textbf{Diagram notation.} Person shape = human actor; cloud = supporting system; blue = selected use case; solid line = association; labelled solid arrow = generalization; dashed arrow = include or extend.
\par\medskip

\hypertarget{sec:a2-actor-register}{}
\addcontentsline{toc}{subsubsection}{A2. Actor Register}
\subsubsection*{A2. Actor Register}
\par\medskip

\begingroup
\small
\renewcommand{\baselinestretch}{1.0}\selectfont
\renewcommand{\arraystretch}{1.15}
\setlength{\tabcolsep}{6pt}
\begin{longtable}{|>{\raggedright\arraybackslash}p{0.09\textwidth}|>{\raggedright\arraybackslash}p{0.09\textwidth}|>{\raggedright\arraybackslash}p{0.74\textwidth}|}
\hline
\textbf{Actor} & \textbf{Type} & \textbf{Description} \\ \hline
\endhead
User & Primary & A registered user who reports inappropriate communication or behavior from profiles, posts, comments, messages, or other users (US-40). \\ \hline
Moderator & Secondary & An authorized reviewer who can review and act on reports according to platform policy (\texttt{AC-40.4}, \texttt{NFR-SEC-10}). \\ \hline
\end{longtable}
\endgroup
\par\medskip

\hypertarget{sec:a3-relationships}{}
\addcontentsline{toc}{subsubsection}{A3. Relationships}
\subsubsection*{A3. Relationships}
\par\medskip

\hypertarget{sec:a3-1-associations}{}
\paragraph*{A3.1 Associations}
\par\medskip

\begingroup
\small
\renewcommand{\baselinestretch}{1.0}\selectfont
\renewcommand{\arraystretch}{1.15}
\setlength{\tabcolsep}{6pt}
\begin{longtable}{|>{\raggedright\arraybackslash}p{0.09\textwidth}|>{\raggedright\arraybackslash}p{0.23\textwidth}|>{\raggedright\arraybackslash}p{0.09\textwidth}|>{\raggedright\arraybackslash}p{0.49\textwidth}|}
\hline
\textbf{Actor} & \textbf{Use Case} & \textbf{Type} & \textbf{Description} \\ \hline
\endhead
User & UC-34 Report Inappropriate Behavior & Primary & Submits a report with category, description, and contextual information. \\ \hline
Moderator & UC-34 Report Inappropriate Behavior & Secondary & Reviews routed reports and may act according to platform policy. \\ \hline
\end{longtable}
\endgroup
\par\medskip

\hypertarget{sec:a3-2-generalization-and-specialization}{}
\paragraph*{A3.2 Generalization and Specialization}
\par\medskip

\begingroup
\small
\renewcommand{\baselinestretch}{1.0}\selectfont
\renewcommand{\arraystretch}{1.15}
\setlength{\tabcolsep}{6pt}
\begin{longtable}{|>{\raggedright\arraybackslash}p{0.09\textwidth}|>{\raggedright\arraybackslash}p{0.09\textwidth}|>{\raggedright\arraybackslash}p{0.09\textwidth}|>{\raggedright\arraybackslash}p{0.65\textwidth}|}
\hline
\textbf{Child} & \textbf{Parent} & \textbf{Type} & \textbf{Inherited Behavior or Association} \\ \hline
\endhead
None & None & None & The proposals do not name an actor generalization or use-case specialization for US-40. \\ \hline
\end{longtable}
\endgroup
\par\medskip

\hypertarget{sec:a3-3-include}{}
\paragraph*{A3.3 Include}
\par\medskip

\begingroup
\small
\renewcommand{\baselinestretch}{1.0}\selectfont
\renewcommand{\arraystretch}{1.15}
\setlength{\tabcolsep}{6pt}
\begin{longtable}{|>{\raggedright\arraybackslash}p{0.13\textwidth}|>{\raggedright\arraybackslash}p{0.16\textwidth}|>{\raggedright\arraybackslash}p{0.63\textwidth}|}
\hline
\textbf{Base Use Case} & \textbf{Included Use Case} & \textbf{Meaning} \\ \hline
\endhead
None & None & The proposals do not name a distinct included use case for US-40. \\ \hline
\end{longtable}
\endgroup
\par\medskip

\hypertarget{sec:a3-4-extend}{}
\paragraph*{A3.4 Extend}
\par\medskip

\begingroup
\small
\renewcommand{\baselinestretch}{1.0}\selectfont
\renewcommand{\arraystretch}{1.15}
\setlength{\tabcolsep}{6pt}
\begin{longtable}{|>{\raggedright\arraybackslash}p{0.15\textwidth}|>{\raggedright\arraybackslash}p{0.11\textwidth}|>{\raggedright\arraybackslash}p{0.10\textwidth}|>{\raggedright\arraybackslash}p{0.56\textwidth}|}
\hline
\textbf{Extending Use Case} & \textbf{Base Use Case} & \textbf{Condition} & \textbf{Added Behavior} \\ \hline
\endhead
None & None & None & The proposals do not name a distinct extending use case for US-40. \\ \hline
\end{longtable}
\endgroup
\par\medskip

\hypertarget{sec:a4-model-notes}{}
\addcontentsline{toc}{subsubsection}{A4. Model Notes}
\subsubsection*{A4. Model Notes}
\par\medskip

\begin{enumerate}
\item UC-34 is derived solely from US-40; no I-xx, E-xx, or G-xx identifier is introduced.
\item The User actor is taken directly from the US-40 story wording "As a user".
\item Moderator is included because \texttt{AC-40.4} explicitly names moderators as reviewers; no separate moderation UC ID is named in the proposals.
\item The moderation queue is a UI surface in \texttt{moderation\_queue.md}, not an external system actor, so it is not drawn as a cloud.
\end{enumerate}
\par\medskip

\hypertarget{sec:part-b-use-case-specification}{}
\addcontentsline{toc}{subsection}{Part B --- Use Case Specification}
\subsection*{Part B --- Use Case Specification}
\par\medskip

\hypertarget{sec:b1-brief-description}{}
\addcontentsline{toc}{subsubsection}{B1. Brief Description}
\subsubsection*{B1. Brief Description}
\par\medskip

The user reports inappropriate communication or behavior so that the platform can maintain a safe and supportive learning environment. US-40 requires reporting from profiles, posts, comments, messages, and other users, categories such as harassment, spam, and inappropriate content, contextual capture, moderation-queue routing, and moderator review under platform policy. The reporter receives confirmation of submission only; no review outcome is promised.
\par\medskip

\hypertarget{sec:b2-preconditions}{}
\addcontentsline{toc}{subsubsection}{B2. Preconditions}
\subsubsection*{B2. Preconditions}
\par\medskip

\begin{enumerate}
\item The user is authenticated with an active session (\texttt{NFR-SEC-1}).
\item The report surface is visible to the user and the target is accessible under privacy and authorization rules.
\item A report action is reachable within at most 2 interactions from a post or comment where applicable (\texttt{NFR-USE-6}).
\item The user can select a supported category and provide optional context.
\end{enumerate}
\par\medskip

\hypertarget{sec:b3-postconditions}{}
\addcontentsline{toc}{subsubsection}{B3. Postconditions}
\subsubsection*{B3. Postconditions}
\par\medskip

\begin{enumerate}
\item \textbf{Success:} The report captures contextual information and is routed to the moderation queue (\texttt{FR-40.3}, \texttt{FR-40.4}, \texttt{AC-40.3}).
\item \textbf{Success:} The reporting user sees confirmation that the report was submitted, without a promised moderator outcome.
\item \textbf{Success:} An authorized moderator can review the report according to platform policy (\texttt{AC-40.4}).
\item \textbf{Failure:} No report is recorded, and the user receives an appropriate error message.
\end{enumerate}
\par\medskip

\hypertarget{sec:b4-main-success-scenario-basic-flow}{}
\addcontentsline{toc}{subsubsection}{B4. Main Success Scenario (Basic Flow)}
\subsubsection*{B4. Main Success Scenario (Basic Flow)}
\par\medskip

\begin{enumerate}
\item The user opens a surface that supports reporting, such as a forum post, comment, profile, message, group chat, or another user.
\item The system displays a report action according to the user's permissions.
\item The user selects \textbf{Report}, \textbf{Report message}, or \textbf{Report inappropriate behavior} as offered on that surface.
\item The system opens \texttt{ReportBehaviorDialog} or \texttt{ReportContentDialog} with a read-only context summary of the target.
\item The user selects one category from \textbf{Spam}, \textbf{harassment}, \textbf{inappropriate content}, or \textbf{academic dishonesty} as listed by \texttt{ReportCategorySelect} in \texttt{moderation\_queue.md}.
\item The user optionally enters a description and submits the report.
\item For a post or comment, the frontend sends \texttt{POST /forum/reports} with \texttt{target\_type}, \texttt{target\_id}, \texttt{category}, and optional \texttt{description}.
\item The backend verifies authentication, target context, category validity, duplicate-report prevention, and reporting rate limits.
\item The backend records the report, returns \texttt{report\_id}, routes it to the moderation queue, shows \texttt{ReportSubmittedNotice} with "Report submitted for review.", and does not state a moderator action.
\item A moderator reviews the report context in the moderation queue and may act according to platform policy.
\end{enumerate}
\par\medskip

\hypertarget{sec:b5-alternative-flows}{}
\addcontentsline{toc}{subsubsection}{B5. Alternative Flows}
\subsubsection*{B5. Alternative Flows}
\par\medskip

\hypertarget{sec:a1-report-a-forum-post-or-comment}{}
\paragraph*{A1. Report a forum post or comment}
\par\medskip

The user selects \texttt{ReportPostButton} on a post or comment. \texttt{ReportContentDialog} collects \texttt{target\_type: 'post' | 'comment'}, \texttt{target\_id}, \texttt{category}, and optional \texttt{description}; this is the only target shape supported by \texttt{POST /forum/reports}.
\par\medskip

\hypertarget{sec:a2-report-a-profile-message-or-other-user}{}
\paragraph*{A2. Report a profile, message, or other user}
\par\medskip

The user selects \texttt{ReportBehaviorButton} from a friend profile, conversation, or member menu. \texttt{ReportBehaviorDialog} captures category, description, and profile or message context; no API endpoint is defined for these targets.
\par\medskip

\hypertarget{sec:a3-report-group-chat-or-member-behavior}{}
\paragraph*{A3. Report group chat or member behavior}
\par\medskip

The user selects \texttt{GroupMessageReportButton}, labelled \textbf{Report message}, from group chat or \texttt{ReportBehaviorButton} from the member list. \texttt{group\_detail.md} routes the report to moderation and shows \texttt{ReportSubmittedNotice} for submission only.
\par\medskip

\hypertarget{sec:b6-exception-flows}{}
\addcontentsline{toc}{subsubsection}{B6. Exception Flows}
\subsubsection*{B6. Exception Flows}
\par\medskip

\hypertarget{sec:e1-invalid-category}{}
\paragraph*{E1. Invalid category}
\par\medskip

\texttt{POST /forum/reports} returns \texttt{400} with \texttt{INVALID\_REPORT\_CATEGORY} (\texttt{48009}). The system keeps the dialog open and asks the user to choose a valid category.
\par\medskip

\hypertarget{sec:e2-target-not-found}{}
\paragraph*{E2. Target not found}
\par\medskip

\texttt{POST /forum/reports} returns \texttt{404} with \texttt{POST\_NOT\_FOUND} (\texttt{48001}) or \texttt{COMMENT\_NOT\_FOUND} (\texttt{48002}). The system explains that the reported content is no longer available.
\par\medskip

\hypertarget{sec:e3-unauthenticated-report-attempt}{}
\paragraph*{E3. Unauthenticated report attempt}
\par\medskip

The system returns \texttt{401} or redirects the user to Login (\texttt{NFR-SEC-1}) and does not submit the report.
\par\medskip

\hypertarget{sec:e4-duplicate-report-or-rate-limit}{}
\paragraph*{E4. Duplicate report or rate limit}
\par\medskip

The backend rejects the report under duplicate-report prevention or \texttt{NFR-SEC-9}. No duplicate-report error code is specified, so the UI shows generic retry-later feedback and does not claim a specific code.
\par\medskip

\hypertarget{sec:e5-server-or-database-failure}{}
\paragraph*{E5. Server or database failure}
\par\medskip

\texttt{POST /forum/reports} returns \texttt{500} with \texttt{DATABASE\_ERROR} or \texttt{UNKNOWN\_ERROR}. The database function also lists \texttt{UNKNOWN\_ERROR} (\texttt{45000}) and \texttt{FORBIDDEN} (\texttt{48010}); no report is confirmed as submitted.
\par\medskip

\hypertarget{sec:b7-business-rules}{}
\addcontentsline{toc}{subsubsection}{B7. Business Rules}
\subsubsection*{B7. Business Rules}
\par\medskip

\begin{enumerate}
\item The system shall allow a user to report inappropriate communication or behavior (\texttt{FR-40.1}, \texttt{AC-40.1}).
\item Reporting is available from forum posts, comments, profiles, and messages (\texttt{FR-40.2}); \texttt{AC-40.1} also names other users.
\item Reports support categories such as harassment, spam, and inappropriate content (\texttt{AC-40.2}); \texttt{moderation\_queue.md} lists \textbf{Spam}, \textbf{harassment}, \textbf{inappropriate content}, and \textbf{academic dishonesty}.
\item The system shall capture contextual information and route reports to the moderation queue (\texttt{FR-40.3}, \texttt{FR-40.4}, \texttt{AC-40.3}).
\item Moderators can review and act on reports according to platform policy (\texttt{AC-40.4}), and reporting supports a safe and supportive university learning environment (\texttt{FR-40.5}).
\item Reporting actions are subject to authentication, server-side authorization, and rate limits (\texttt{NFR-SEC-1}, \texttt{NFR-SEC-6}, \texttt{NFR-SEC-9}, \texttt{NFR-SEC-10}).
\item Report submission feedback must not promise a moderator outcome, and privacy settings apply consistently to report context (\texttt{NFR-PRIV-7}).
\end{enumerate}
\par\medskip

\hypertarget{sec:b8-special-requirements}{}
\addcontentsline{toc}{subsubsection}{B8. Special Requirements}
\subsubsection*{B8. Special Requirements}
\par\medskip

\begingroup
\small
\renewcommand{\baselinestretch}{1.0}\selectfont
\renewcommand{\arraystretch}{1.15}
\setlength{\tabcolsep}{6pt}
\begin{longtable}{|>{\raggedright\arraybackslash}p{0.10\textwidth}|>{\raggedright\arraybackslash}p{0.83\textwidth}|}
\hline
\textbf{ID} & \textbf{Requirement} \\ \hline
\endhead
NFR-SEC-1, NFR-SEC-3 & Authentication is required before protected report surfaces are accessible; sensitive communication is encrypted using TLS 1.2 or higher. \\ \hline
NFR-SEC-6 & Unauthorized access is prevented through server-side authorization checks. \\ \hline
NFR-SEC-9 & Reporting and related actions are rate-limited to reduce spam and abuse. \\ \hline
NFR-SEC-10 & Moderation actions are restricted to authorized moderators and administrators. \\ \hline
NFR-USE-6, NFR-PRIV-7 & Reporting actions are reachable within at most 2 interactions from a post or comment, and privacy settings apply consistently to forum posts, comments, member lists, and search results. \\ \hline
Design constraint & Report UI confirms submission only and must not state or imply a moderator decision. \\ \hline
\end{longtable}
\endgroup
\par\medskip

\hypertarget{sec:b9-extension-points}{}
\addcontentsline{toc}{subsubsection}{B9. Extension Points}
\subsubsection*{B9. Extension Points}
\par\medskip

\begingroup
\small
\renewcommand{\baselinestretch}{1.0}\selectfont
\renewcommand{\arraystretch}{1.15}
\setlength{\tabcolsep}{6pt}
\begin{longtable}{|>{\raggedright\arraybackslash}p{0.37\textwidth}|>{\raggedright\arraybackslash}p{0.56\textwidth}|}
\hline
\textbf{Extension Point} & \textbf{Location in Flow of Events} \\ \hline
\endhead
Forum post or comment entry & Basic flow steps 1-7; alternative A1 \\ \hline
Profile, message, or user entry & Basic flow steps 1-6; alternative A2 \\ \hline
Group chat or member entry & Basic flow steps 1-6; alternative A3 \\ \hline
Category and target validation & Basic flow steps 5 and 7; exceptions E1 and E2 \\ \hline
\end{longtable}
\endgroup
\par\medskip

\hypertarget{sec:b10-interface-and-data-notes}{}
\addcontentsline{toc}{subsubsection}{B10. Interface and Data Notes}
\subsubsection*{B10. Interface and Data Notes}
\par\medskip

\begin{itemize}
\item \textbf{Screen, flow, and entry points.} \texttt{screen\_flow.md} places report-capable content under Home $\rightarrow$ Forum Group and Home $\rightarrow$ Friend Group; \texttt{feature\_flow.md} gives \texttt{Post -$>$ Report / Moderate}; \texttt{userstory.md} section 11 places Notification Center activity beside forum and social activity. \texttt{settings.md} defines \texttt{ReportBehaviorLink} with label \textbf{Report inappropriate behavior}; \texttt{friend.md} and \texttt{friend\_profile.md} define \texttt{ReportBehaviorButton} with label \textbf{Report}; \texttt{group\_detail.md} defines \texttt{GroupMessageReportButton} with label \textbf{Report message} and \texttt{ReportBehaviorDialog}.
\item \textbf{Exact categories.} \texttt{moderation\_queue.md} \texttt{ReportCategorySelect} lists \textbf{Spam}, \textbf{harassment}, \textbf{inappropriate content}, and \textbf{academic dishonesty}. \texttt{AC-40.2} names harassment, spam, and inappropriate content; \texttt{FR-22.2}/\texttt{AC-22.1} also name academic dishonesty.
\item \textbf{Report API.} \texttt{POST /forum/reports} body is \texttt{\{target\_type: 'post' | 'comment', target\_id: string, category: string, description?: string\}}; success is \texttt{201 \{report\_id\}}. Errors are \texttt{400 INVALID\_REPORT\_CATEGORY} (\texttt{48009}), \texttt{404 POST\_NOT\_FOUND} (\texttt{48001}) or \texttt{COMMENT\_NOT\_FOUND} (\texttt{48002}), \texttt{401}, and \texttt{500 (DATABASE\_ERROR | UNKNOWN\_ERROR)}.
\item \textbf{Database contract and data fields.} \texttt{studitopia.report\_forum\_content(preporter\_id, ptarget\_type, ptarget\_id, pcategory, pdescription)} returns \texttt{\{report\_id: uuid\}} and lists \texttt{UNKNOWN\_ERROR} (\texttt{45000}), \texttt{POST\_NOT\_FOUND} (\texttt{48001}), \texttt{COMMENT\_NOT\_FOUND} (\texttt{48002}), \texttt{INVALID\_REPORT\_CATEGORY} (\texttt{48009}), and \texttt{FORBIDDEN} (\texttt{48010}). Duplicate-report prevention is required but has no error code. \texttt{moderation\_queue.md} lists \texttt{report\_id}, \texttt{target\_type}, \texttt{target\_id}, \texttt{category}, \texttt{description}, \texttt{reporter\_id}, \texttt{created\_at}, \texttt{status}, \texttt{moderation\_action}, \texttt{action\_reason}, and \texttt{affected\_user\_notified}; \texttt{status}, queue, action, and history contracts are gaps.
\item \textbf{Submission feedback.} \texttt{ReportSubmittedNotice} says "Report submitted for review." (\texttt{friend\_profile.md}, \texttt{group\_detail.md}). This confirms submission only; reviewer actions are outside the reporter UI.
\item \textbf{API gaps and recommendations.} Profiles, messages, other users, and group activity have report requirements but no report endpoint in \texttt{API.md}; no queue, action, history, or status endpoint is defined. Do not invent an endpoint, status code, or error code for those surfaces. \texttt{/moderation/queue}, \texttt{ModerationQueuePage}, \texttt{ModerationQueueTable}, \texttt{ReportDetailPanel}, and \texttt{ModerationActionDialog} are recommendations in \texttt{moderation\_queue.md}. Moderation actions require \texttt{NFR-SEC-10}, server-side checks require \texttt{NFR-SEC-6}, and report context must respect \texttt{NFR-PRIV-7}.
\end{itemize}
\par\medskip

\hypertarget{sec:b11-source-alignment-notes}{}
\addcontentsline{toc}{subsubsection}{B11. Source Alignment Notes}
\subsubsection*{B11. Source Alignment Notes}
\par\medskip

\begin{itemize}
\item \texttt{userstory.md} section 11.3 and \texttt{functional\_requirement.md} US-40 agree on report surfaces, categories, context capture, routing, and moderator review; \texttt{moderation\_queue.md} adds the only defined report API contract and also cites \texttt{FR-22.1}-\texttt{FR-22.6} and \texttt{AC-22.1}-\texttt{AC-22.5}, overlapping the report API notes in existing PA4 UC-05.
\item No communication-settings or behavioral-report API endpoints are defined for this use case: \texttt{API.md} has no communication-settings endpoint and no behavioral-report endpoint for profiles, messages, other users, or group activity; only the forum-specific \texttt{POST /forum/reports} exists for posts and comments.
\item Report UI must not promise a moderator outcome; \texttt{ReportSubmittedNotice} gives submission confirmation only, and \texttt{AC-40.4} leaves the decision to authorized moderators under platform policy.
\item Forum-specific content reports overlap with PA4 \texttt{UC-20-report-moderate-forum-content.md}, which was not modified; this use case covers behavioral reports from profiles, messages, groups, and other users under US-40, and \texttt{friend.md}, \texttt{friend\_profile.md}, and \texttt{group\_detail.md} are UI/UX recommendations with no friend/profile/group API or error-code contract. \texttt{moderation\_queue.md} states duplicate-report prevention without an error code, so generic failure is used here.
\item No distinct include, extend, or specialization use case is named for US-40, so no I-xx, E-xx, or G-xx identifier is created.
\end{itemize}
\par\medskip

\par\medskip
\hypertarget{sec:ui-prototype-specimens}{}
\addcontentsline{toc}{subsection}{UI Prototype Specimens}
\subsection*{UI Prototype Specimens}
\par\medskip
\begin{figure}[H]
\centering
\includegraphics[width=0.75\textwidth,height=0.42\textheight,keepaspectratio]{"../../prototype/profile_p5_report_a_user.png"}
\par\vspace{2pt}
{\footnotesize \textbf{UI Prototype:} Profile: P5 Report a user}
\end{figure}
\par\smallskip

\end{document}
